\bookchapter{บทที่ 19: การไล่หาปัญหาและช่วยเหลือตัวเอง}

\chapterintro{มือถือก็เหมือนคน เป็น “ป่วย” ได้เหมือนกัน กระตุก แอปปิดตัว จอดำ ชาร์จไม่เข้า… ปัญหาส่วนใหญ่ไม่น่ากลัวอย่างที่คิด ก่อนส่งซ่อมคุณลองเองก่อนได้ทั้งนั้น บทนี้คือ “คู่มือฉุกเฉินมือถือ” ของคุณ — ลองวิธีต้นทุนต่ำก่อน ไม่ได้ค่อยหาหมอมืออาชีพ}

\booksection{หนึ่ง มือถือกระตุก ร้อน — คำถามที่ถูกถามบ่อยสุด}

พูดตรงๆ มือถือใช้ไปราวหนึ่งปีเริ่มกระตุก เป็นเรื่องปกติมาก สาเหตุไม่จำเป็นต้องเป็นเครื่องพัง มักเป็นการ “ฝุ่นสะสม” — ไฟล์ขยะ โปรเซสพื้นหลัง พื้นที่จัดเก็บค่อยๆ เต็ม

\booksubsection{วินิจฉัยก่อน อย่ารีบซ่อม:}

\begin{enumerate}
  \item \textbf{ตรวจพื้นที่จัดเก็บ}: ไปที่ตั้งค่า พื้นที่จัดเก็บ ดูพื้นที่เหลือ ถ้าเหลือไม่ถึง 10\% ระบบจะกระตุกมาก เพราะมันกำลังดิ้นรนหาพื้นที่ ประสบการณ์ของฉัน: รักษาพื้นที่ว่างไว้อย่างน้อย 15\%-20\% เครื่องถึงจะลื่นพอสมควร
  \item \textbf{ดูอันดับการกินไฟ}: ไปที่การตั้งค่าแบตเตอรี่ดูอันดับการกินไฟ 24 ชั่วโมงที่ผ่านมา ถ้าแอปไหนกินเป็นหัวคะแนนแต่คุณแทบไม่ได้ใช้ แปลว่ามันแอบวิ่งในพื้นหลัง ถอนการติดตั้งหรือจำกัดกิจกรรมพื้นหลังเลย
  \item \textbf{นึกดูว่าเพิ่งทำอะไรมา}: อัปเดตระบบแล้วกระตุก? (หลังอัปเดตเวอร์ชันใหญ่กระตุกสองสามวันแรกเป็นปกติ ระบบกำลังสร้างดัชนีใหม่ในพื้นหลัง) หรือติดตั้งแอปใหม่แล้วกระตุก? (ลองถอนดู)
  \item \textbf{ดูสภาพแวดล้อมการใช้งาน}: มือถือตากแดดจัด ใส่เคสหนาชาร์จไปเล่นเกมไป ร้อนเป็นธรรมดา ความร้อนยังทำให้ CPU ลดความถี่อัตโนมัติ คุณจะรู้สึกว่าเครื่องกระตุกฉับพลัน พอเย็นลงค่อยลองใหม่
\end{enumerate}

\textbf{วิธีรักษา (จากง่ายไปยาก)}: รีบูตเครื่อง (แก้ปัญหาได้ 60\%) เคลียร์พื้นที่จัดเก็บ (ลบแอปกับวิดีโอที่ไม่ใช้) แอปไหนสร้างปัญหาก็ล้างแคชมันหรือถอนติดตั้งใหม่ ท่าสุดท้าย: สำรองข้อมูลแล้วรีเซ็ตโรงงาน

\booksection{สอง แอปปิดตัวเอง ไม่ตอบสนอง — แอปไม่ฟังคำสั่ง}

แอปไหนเปิดปุ๊บปิดปั๊บ หรือใช้ไปใช้มาค้างแน่นิ่ง กระบวนการจัดการเหมือนหมอซักประวัติ ทำทีละขั้น:

\begin{enumerate}
  \item \textbf{ปิดแล้วเปิดใหม่ก่อน}: ปัดแอปออกจากพื้นหลังแล้วเปิดใหม่ บั๊กชั่วคราวส่วนใหญ่จบด้วยท่านี้
  \item \textbf{อัปเดตแอป}: ไปสโตร์ดูว่ามีอัปเดตไหม บางทีเป็นเพราะเวอร์ชันแอปเก่าเกินไป ไม่เข้ากับระบบ
  \item \textbf{ตรวจสิทธิ์}: แอปบางตัวต้องมีสิทธิ์ใดสิทธิ์หนึ่งถึงทำงานปกติ แต่คุณไม่ได้ให้ เช่นแอปกล้องไม่ได้รับสิทธิ์กล้อง เปิดปุ๊บปิดปั๊บ ไปตั้งค่าหาแอปนั้น ให้สิทธิ์ที่ต้องให้ให้ครบ
  \item \textbf{ล้างแคช}: ไปที่ ตั้งค่า จัดการแอป หาแอปนั้น พื้นที่จัดเก็บ ล้างแคช ระวังคือ "ล้างแคช" ไม่ใช่ "ล้างข้อมูล"! ล้างข้อมูลจะลบสถานะล็อกอินกับการตั้งค่าในแอปทั้งหมด เท่ากับรีเซ็ตโรงงานเฉพาะแอป
  \item \textbf{ถอนแล้วติดตั้งใหม่}: ท่าสุดท้าย ก่อนถอนยืนยันก่อนว่าข้อมูลในแอปจะหายไหม — WeChat สำรองประวัติแชทก่อน เกมยืนยันก่อนว่าผูกบัญชีกับเบอร์โทรแล้ว
\end{enumerate}

\booksection{สาม ชาร์จไม่เข้าหรือชาร์จช้า — อาจไม่ใช่ความผิดแบต}

ชาร์จไม่เข้าเป็นปัญหาที่ทำให้ใจร้อนง่ายมาก แต่หลายครั้งสาเหตุอยู่ที่อุปกรณ์เสริม ไม่ใช่เครื่องพัง

\begin{itemize}
  \item \textbf{ตรวจช่องชาร์จก่อน}: ใช้ไม้จิ้มฟันหรือเครื่องมือเล็กที่ไม่นำไฟ แซะฝุ่นกับปุยฝ้ายในช่องชาร์จเบาๆ อันนี้ฉันเคยเจอเอง — เศษด้ายในกระเป๋าตกไปอุดช่องชาร์จพอดี เสียบเข้าไปคิดว่าขั้วต่อไม่ดี แซะออกก็หาย
  \item \textbf{เปลี่ยนสายกับหัวชาร์จลองดู}: สายเสียหายง่ายกว่าหัวชาร์จมาก โดยเฉพาะสายที่ม้วนบิดไปมาบ่อย เปลี่ยนสายกับหัวชาร์จที่แน่ใจว่าปกติมาทดสอบไขว้ เร็วๆ นี้ก็รู้ว่าอุปกรณ์ไหนพัง
  \item \textbf{ปัญหาการชาร์จไร้สาย}: ถ้าใช้ชาร์จไร้สาย ตรวจว่าเคสหนาเกินไหม หรือหลังเคสมีแผ่นโลหะ (ขาแม่เหล็กพวกนี้) ไหม และวางมือถือให้ตรงกลางแผ่นชาร์จ — เอียงไปอาจชาร์จไม่เข้า
  \item \textbf{ชาร์จช้า}: ยืนยันว่ากำลังหัวชาร์จพอไหม (iPhone จะชาร์จเร็วต้องใช้หัว PD อย่างน้อย 20W) สายรองรับชาร์จเร็วไหม (ไม่ใช่ทุกสายจะทนกระแนเกิน 3A) มือถือร้อนเกินไปไหม (ร้อนเกินจะลดกำลังอัตโนมัติเพื่อปกป้องแบต)
\end{itemize}

\booksection{สี่ เน็ตต่อไม่ได้ Wi-Fi หลุดบ่อย}

นอกบ้านใช้ดาต้าได้ปกติ กลับถึงบ้านต่อ Wi-Fi ไม่ได้ หรือ Wi-Fi เชื่อมอยู่แต่ไม่มีเน็ต ไล่ตรวจตามลำดับนี้:

\begin{enumerate}
  \item \textbf{เปิดปิดโหมดเครื่องบิน}: ปัดลงศูนย์ควบคุม เปิดโหมดเครื่องบินรอ 5 วินาทีแล้วปิด เท่ากับให้มือถือค้นหาเครือข่ายใหม่อีกรอบ ได้ผลมาก
  \item \textbf{Wi-Fi "ลืม" แล้วเชื่อมใหม่}: ไปที่การตั้งค่า Wi-Fi กดเครื่องหมายตกใจข้างชื่อเครือข่าย (แอนดรอยด์คือไอคอนฟันเฟือง) เลือก "ละเว้นเครือข่ายนี้" หรือ "ลืม" แล้วกรอกรหัสผ่านเชื่อมใหม่
  \item \textbf{รีบูตเราเตอร์}: หลายบ้านเราเตอร์ไม่ได้ปิดเป็นเดือน นานเข้ามันจะไม่ค่อยนิ่ง ถอดปลั๊กรอ 30 วินาทีแล้วเสียบกลับ ปัญหา Wi-Fi ส่วนใหญ่จบ
  \item \textbf{ตรวจว่าเป็นที่มือถือจริงไหม}: อุปกรณ์อื่นต่อ Wi-Fi เดียวกันปกติไหม? ถ้าปกติทั้งหมดแต่มือถือคุณไม่ได้ เครื่องอาจต้องรีบูตหรือรีเซ็ตการตั้งค่าเครือข่าย (ตั้งค่า ทั่วไป โอนย้ายหรือรีเซ็ต รีเซ็ตการตั้งค่าเครือข่าย ไม่ลบข้อมูลของคุณ แต่จะล้างรหัส Wi-Fi ทั้งหมด)
\end{enumerate}

\booksection{ห้า ปัญหาหน้าจอ: สัมผัสไม่ไว กะพริบ มีเงาติด}

\begin{itemize}
  \item \textbf{สัมผัสไม่ค่อยไว}: ลองแกะฟิล์มกระจกออกดูว่าเป็นที่ฟิล์มไหม — ฟิล์มถูกๆ ใช้นานๆ กระทบการสัมผัสจริงๆ หน้าจอมีน้ำหรือมือมันก็ทำให้สัมผัสเพี้ยนได้ ไม่ใช่ทั้งคู่ก็รีบูตแล้วลองใหม่
  \item \textbf{จอมีเงา/ภาพค้าง (จอไหม้)}: แสดงภาพเดิมนานๆ (เช่นเปิดจอปัด TikTok ตลอด) จะทิ้ง "ภาพค้าง" บนจอ OLED นี่ไม่ใช่จอพัง เป็นแค่ชั่วคราว ปิดจอพักครึ่งชั่วโมงมักหาย
  \item \textbf{หน้าจอกะพริบ}: ถ้ากะพริบตอนความสว่างต่ำ อาจเป็นปัญหาการปรับแสง PWM พบมากบนจอ OLED ไปตั้งค่าหา "DC dimming" หรือ "กันกระพริบ" เปิดแล้วดีขึ้นเยอะ ไม่ใช่ทุกเครื่องมีฟังก์ชันนี้
  \item \textbf{หน้าจอมีบล็อกหนึ่งดำสนิท/สว่างผิดปกติ}: อันนี้หนักข้อ น่าจะจอเสียหายทางกายภาพ ต้องเปลี่ยนจอ จอเปลี่ยนนอกทางการสีและความสว่างมักสู้ของแท้ไม่ได้ แถมกระทบความกันน้ำ
\end{itemize}

\booksection{หก เปิดเครื่องไม่ได้ ค้างที่โลโก้ — อย่าเพิ่งตกใจ}

สถานการณ์นี้ทำให้หัวใจเต้นแรงที่สุด แต่ใจเย็น ลองพวกนี้ก่อน:

\begin{itemize}
  \item \textbf{ชาร์จครึ่งชั่วโมงก่อนค่อยว่ากัน}: หลายครั้งแค่แบตหมดเกลี้ยง เสียบหัวชาร์จของแท้ รออย่างน้อย 30 นาที แล้วค่อยลองเปิดเครื่อง บางเครื่องพอแบตหมดต้องชาร์จถึงระดับหนึ่งถึงจะตอบสนอง
  \item \textbf{รีบูตแบบบังคับ}: iPhone จอเต็มกดเพิ่มเสียงหนึ่งครั้ง ลดเสียงหนึ่งครั้ง กดปุ่มพาวเวอร์ด้านข้างค้างไม่ปล่อย จนจอโผล่โลโก้แอปเปิล (อาจต้องกดค้าง 10-15 วินาที) วิธีทั่วไปแอนดรอยด์: กดปุ่มพาวเวอร์ค้างเกิน 10 วินาที จนสั่นหรือจอสว่าง
  \item \textbf{ค้างหน้าเปิดเครื่องรีบูตวนไปมา}: สถานการณ์นี้เรียก "ขาวแอปเปิล" (iPhone) หรือ "ค้างโลโก้" (แอนดรอยด์) ถ้ารีบูตบังคับไม่หาย อาจต้องเข้าโหมดกู้คืนซ่อมด้วยคอม iPhone ต่อคอมเปิด iTunes/Finder เข้าโหมดกู้คืนแล้วเลือก "อัปเดต" (ไม่ใช่ "กู้คืน") จะติดตั้งระบบใหม่แต่เก็บข้อมูลไว้ แอนดรอยด์เข้า Fastboot หรือ Recovery คีย์ลัดแต่ละแบรนด์ต่างกัน แนะนำค้นหารุ่นมือถือคุณ + "เข้าโหมด Recovery"
  \item \textbf{ลองหมดแล้วไม่ได้}: ก็ต้องขอความช่วยเหลือจากร้านซ่อมมืออาชีพ ถ้าสำรองข้อมูลไว้แล้วก็ไม่ต้องกังวลมาก
\end{itemize}

\booksection{เจ็ด ปัญหาเสียง: ลำโพงเงียบ เสียงโทรเบา}

\begin{itemize}
  \item \textbf{ลำโพงไม่มีเสียง}: ยืนยันก่อนว่าไม่ใช่โหมดเงียบ (สวิตช์เงียบข้างเครื่อง iPhone) โหมดห้ามรบกวน หรือเสียงถูกกดต่ำสุด ฟังดูโง่ แต่กรณีพวกนี้คิดเป็น 40\% แล้วตรวจว่าบลูทูธกำลังเชื่อมหูฟังหรือลำโพงตัวไหนอยู่
  \item \textbf{โทรแล้วอีกฝ่ายไม่ได้ยินคุณ}: ตรวจว่ารูไมค์ถูกปิดหรือเปล่า — รูเล็กๆ ด้านล่างเครื่องถูกฝุ่นปิดง่ายมาก ใช้แปรงขนนุ่มเก็บฝุ่นเบาๆ อย่าใช้เข็มหรือไม้จิ้มฟันเอาแรงแทง จะทะลุตาข่ายกันฝุ่น
  \item \textbf{เสียงแหบ/แตก}: ถ้าลำโพงโดนน้ำ — อย่าใช้ไดรเยอร์ลมร้อนเป่า! ใช้กระดาษทิชชูซับน้ำผิวๆ แล้ววางปล่อยแห้งเอง บนเน็ตว่าแอปไล่น้ำด้วยคลื่นเสียงได้ผลพอควร (เล่นเสียงความถี่ต่ำสั่นน้ำออก) ลองได้
\end{itemize}

\booksection{แปด หาซ่อมทางการ เลี่ยงกับดักพวกนั้น}

ถ้าวิธีข้างบนลองหมดแล้วไม่หาย ก็ต้องหาร้านซ่อมมืออาชีพ แต่วงการซ่อมน้ำลึก เล่ากับดักและประสบการณ์ที่ฉันเคยโดน:

\begin{itemize}
  \item \textbf{หาตัวจริงทางการ อย่าหา "หน้าตาเหมือนทาการ"}: ไปเว็บทางการของแบรนด์ (Apple.com.cn, Huawei.com.cn ฯลฯ) หารายชื่อ "ผู้ให้บริการที่ได้รับอนุญาต" ร้านข้างทางแขวนโลโก้แอปเปิลที่ไม่ได้รับอนุญาต อุปกรณ์ที่ใช้ก็อาจไม่ใช่ของแท้ ซ่อมเสร็จความกันน้ำของคุณก็พังไปหนึ่งเส้น
  \item \textbf{สิ่งที่ต้องทำก่อนส่งซ่อม}: สำรองข้อมูลทั้งหมด (อย่าหวังว่าตอนซ่อมข้อมูลจะไม่หาย) ล็อกเอาต์บัญชีสำคัญทั้งหมด (Apple ID, WeChat, Alipay) ปิด "ค้นหา iPhone" บางร้านจะขอรหัสล็อกหน้าจอ เรื่องนี้มีความเสี่ยงด้านความเป็นส่วนตัว — ควรล็อกเอาต์บัญชีเองก่อนแล้วค่อยส่งไป
  \item \textbf{กับดักการเปลี่ยนหน้าจอ}: จอนอกทางการแบ่ง "จอแท้เปลี่ยนแก้ว" (แก้วเดิมเปลี่ยน) "จอเปลี่ยนแก้ว" (แก้วไม่แท้) "จอประกอบ" (ไม่แท้ทั้งชุด) ราคาต่างกันหลายเท่า คุณภาพการแสดงผลกับความทนก็ต่างกันมาก ถ้าเครื่องยังอยู่ในประกัน ซื่อสัตย์เดินทางการเถอะ
  \item \textbf{กับดักการเปลี่ยนแบต}: แบตไม่แท้อาจป้ายความจุเกินจริง ใช้ครึ่งปีเสื่อมกว่าแบตแท้สองปีซะอีก ราคาเปลี่ยนแบตทางการแอปเปิลแพงกว่านอกทางการไม่มาก แนะนำเดินทางการ
  \item \textbf{กับดักการกู้ข้อมูล}: มือถือโดนน้ำ/ตกพื้นแล้วไปซ่อม ถ้าอีกฝ่ายบอก "กู้ข้อมูลต้องจ่ายเพิ่ม" ระวัง — ร้านไม่ซื่อบางร้านขยายความเสียหายเกินจริงเพื่อเก็บเงินเพิ่ม การกู้ข้อมูลที่ถูกต้องคือให้ราคาก่อนแล้วค่อยลงมือ และกู้ไม่สำเร็จไม่คิดเงิน
\end{itemize}

\begin{tcbwarn}{ท่าแรกสารพัดประโยชน์: รีบูตเครื่อง}
ไม่ว่าเจอปัญหาอะไร — กระตุก แอปปิดตัว ต่อเน็ตไม่ได้ เสียงผิดปกติ — รีบูตมือถือก่อน ประโยคนี้ฟังดูแก้ตัว แต่ตามประสบการณ์ของฉันเอง อย่างน้อยหกทบสิบของปัญหาเล็กๆ รีบูตแล้วหาย มือถือเหมือนคอม วิ่งนานๆ หน่วยความจำจะมีขยะสะสมเต็มไปหมด รีบูตคือการเคลียร์เริ่มใหม่ ดังนั้นครั้งหน้ามือถือมีปัญหา อย่ารีบค้น “ทำยังไงดี” รีบูตจบแล้วค่อยดู

\end{tcbwarn}

\begin{tcbpractice}{จดไว้เลยตอนนี้: การช่วยชีวิตพวกนี้ต้องจำขึ้นใจ}
\begin{enumerate}
  \item \textbf{รีบูตแบบบังคับ}: iPhone จอเต็มกดเพิ่มเสียง ลดเสียง กดปุ่มพาวเวอร์ค้างจนโลโก้ขึ้น แอนดรอยด์กดปุ่มพาวเวอร์ค้าง 10-15 วินาทีจนสั่น
  \item \textbf{เก็บฝุ่นช่องชาร์จ}: หาเครื่องมือเล็กๆ (ไม้จิ้มฟันหรือเข็มดึงซิม) เก็บฝุ่นในช่องชาร์จเบาๆ
  \item \textbf{เชื่อม Wi-Fi ใหม่}: ไปที่การตั้งค่า Wi-Fi "ลืม" Wi-Fi ปัจจุบัน แล้วกรอกรหัสผ่านเชื่อมใหม่
  \item \textbf{เปิด/ปิดโหมดเครื่องบิน}: ปัดลงศูนย์ควบคุม เปิดโหมดเครื่องบิน นับในใจ 5 วินาที ปิด ดูว่าเน็ตกลับมาไหม
  \item สกรีนช็อตขั้นตอนการทำงานพวกนี้เก็บไว้ หรือจดไว้ในแอปโน้ต เจอปัญหาไม่ต้องตื่นตระหนกไปค้นหา เปิดโน้ตดูก็ได้
\end{enumerate}

\end{tcbpractice}
